\chapter{Introduction}%
\label{chap:intro}

\emph{Interactive theorem provers} (ITPs) are software systems for rigorous mathematical proof.
They allow their users to express mathematical statements and proofs in a precise, unambiguous language, and check these proofs according to the rules of the respective prover's underlying logic.
The logic is usually kept simple, admitting a small implementation of a proof checker, called a \emph{kernel}.
Hence, if an ITP accepts a proof, we can be very confident that the proof is, in fact, correct.
Indeed, while bugs have been found in many kernels (usually small and quickly fixed), to my knowledge no ITP-proved statement of wider interest has ever turned out to be false.
The same cannot be said for regular, non-computerised mathematics~\cite{Voevodsky14,Hales08} or for many of the automated systems that aim to verify various properties of computer programs~\cite{CSmith}. % TODO better refs

Such extreme reliability has made ITPs a mainstay of certain scientific areas, particularly programming languages research, and they have also found some success in industrial applications where correctness is paramount~\cite{seL4,CompCert,AWSFormal}.
A promising line of research has recently started to combine ITPs with machine learning, aiming to mitigate the unreliability of large language models by checking their outputs with ITPs~\cite{Lample22,LeanDojo,GoedelProver,DeepSeekProver,Aristotle}.
The mathematics community has historically, by and large, been sceptical of computer-based proofs, despite some early successes~\cite{FourColor,Kepler,OddOrder}, but there has recently been renewed interest in ITPs and particularly in the Mathlib project~\cite{Mathlib}.
This mathematics library for the Lean~4 ITP~\cite{Lean4} is perhaps the first large, coherent library of mathematics formalised in an interactive theorem prover and primarily developed by mathematicians rather than computer scientists.
It aims to develop a typical undergraduate mathematics curriculum as well as more advanced concepts, to provide a foundation for both formal developments of research mathematics~\cite{LiquidTensor,PFR} and future AI-based mathematics tools.

ITPs are also called \emph{proof assistants}, hinting at a vision that goes beyond reliability.
Compared to pen-and-paper mathematics, proof assistants can provide more immediate feedback to their users, indicating when and why an attempted proof step is wrong.
They can also help users track the context and structure of their proofs, e.g.\ which cases need to be considered when a case distinction is made or which hypotheses are available at a given point in the proof.

However, this vision of a helpful mathematical assistant remains only partly realised.
Extensive experience is still required for even a trained mathematician to become a productive user of ITPs.
The main cause of this is perhaps the requirement to spell out proofs in full detail, leaving nothing as an exercise for the reader.
Compounding this issue, the minimalistic logics of proof assistants typically require explicit proof steps even for operations that are entirely obvious to humans.
For example, if, for two natural numbers $n$ and $m$, the sum $m + n$ has been shown to be positive, most proof assistants require an explicit proof step to show that $n + m$ is also positive.
Similarly, if two functions $f$ and $g$ are continuous, multiple proof steps are needed to show that $x ↦ f(2g(x) + 1)$ is also continuous.
Demonstrating such trivialities takes time and effort, clutters the proof and saps the enthusiasm of aspiring ITP users.

Almost all proof assistants address this issue with some form of proof automation.
This means the proof assistants offer programs, usually called \emph{tactics}, that partially or fully automate the construction of proofs.
A simple tactic might derive $n + m > 0$ from $m + n > 0$, constructing a proof that uses commutativity of addition and a congruence rule.
A more complex tactic might prove continuity of $x ↦ f(2g(x) + 1)$ by referring to a lemma about the composition of continuous functions, together with continuity of multiplication and addition.

\section{Aesop}%
\label{sec:intro:aesop}

In this thesis, we introduce \emph{Aesop}\footnote{\url{https://github.com/leanprover-community/aesop}} (Automatic Extensible Search for Obvious Proofs), a tactic for the Lean~4 proof assistant.
As the name suggests, Aesop's purpose is to automate the mathematically trivial goals that arise in many proofs -- it is not expected to prove interesting theorems on its own.

When searching for a proof, Aesop systematically applies \emph{rules} that reduce a \emph{goal} to zero or more subgoals.
Each goal consists of a statement to be proved (the \emph{target}) and a list of \emph{hypotheses} that are assumed to be true (the \emph{local context}).
We write a goal with local context $Γ$ and target $T$ as $Γ ⊢ T$.
In their most general form, rules are arbitrary tactics, but Aesop also provides \emph{rule builders} that create rules for specific lemmas.
Rules are added by users, who often curate \emph{rule sets} for certain tasks.

For example, the lemma that continuity of $f$ and $g$ implies continuity of $x ↦ f(x) + g(x)$ (\textsc{Add}) can be turned into a rule that reduces any goal of the form \enquote{$x ↦ f(x) + g(x)$ is continuous} into two subgoals, \enquote{$f$ is continuous} and \enquote{$g$ is continuous}.
Similarly, we can register the following lemmas about continuity:
\begin{itemize}
  \item \textsc{Mul}: If $f$ and $g$ are continuous, then $x ↦ f(x)g(x)$ is continuous.
  \item \textsc{Comp}: If $f$ and $g$ are continuous, then $x ↦ f(g(x))$ is continuous.
  \item \textsc{Const}: The constant function $x ↦ c$ is continuous for any $c$.
\end{itemize}
By repeatedly applying these rules, Aesop can prove that $x ↦ f(f(x)g(x) + 1)$ is continuous.

\begin{figure}
  \centering
  \begin{tikzpicture}[outer sep=auto, level distance=10mm, sibling distance=15mm]
    \node {$x ↦ f(f(x)g(x) + 1)$ is continuous}
    child {node[rectangle,draw] {\textsc{Comp}}
      child[xshift=-3cm] {node {$f$ is continuous}
        child {node[rectangle,draw] {assumption}}}
      child[xshift=3cm] {node {$x ↦ f(x)g(x) + 1$ is continuous}
        child {node[rectangle,draw] {\textsc{Add}}
          child[xshift=-2cm] {node {$x ↦ f(x)g(x)$ is continuous}
            child {node[rectangle,draw] {\textsc{Mul}}
              child[xshift=-1.5cm] {node {$f$ is continuous}
                child {node[rectangle,draw] {assumption}}}
              child[xshift=1.5cm] {node {$g$ is continuous}
                child {node[rectangle,draw] {assumption}}}}}
          child[xshift=2cm] {node {$x ↦ 1$ is continuous}
            child {node[rectangle,draw] {\textsc{Const}}}}}}};
  \end{tikzpicture}
  \caption{Proof of continuity of a particular function}%
  \label{fig:intro:proof-tree}
\end{figure}

We can visualise the proof as the tree shown in Fig.~\ref{fig:intro:proof-tree}, which contains two alternating kinds of nodes, \emph{goal nodes} and \emph{rule application nodes} (also called \emph{rapp nodes}).
A goal node represents a goal; the root goal is our initial conjecture.
A rapp node represents a rule that was applied to prove its parent goal.
The rule may require additional goals to be solved, which become children of the rapp node.
The leaf rapp nodes have no children, so nothing remains to be proved, and the tree represents a complete proof of the initial conjecture.

Usually, Aesop is provided not with exactly the rules needed for a proof, but rather with a superset designed to prove a large class of goals.
As a result, multiple rules may apply to any given goal.
Aesop then considers all possible rule applications as sibling rapp nodes, and the proof tree becomes a search tree.
More specifically, it becomes an AND-OR-tree: sibling goals are implicitly conjoined (all need to be proved for the parent rapp to be proved) while sibling rapps are implicitly disjoint (only one needs to be proved for the parent goal to be proved).

Essentially the same setup is used by various other proof search tactics, particularly Rocq's~\cite{Rocq} and Isabelle's~\cite{Isabelle} auto tactics~\cite{IsabelleAuto}.
Aesop distinguishes itself by three features, summarised in the next three sections and detailed in Chapters~\ref{chap:aesop}--\ref{chap:script}.

\section{Best-First Tree Search with Metavariables}

When searching for proofs, we need to answer two prioritisation questions: which rule do we apply first to a given goal, and which goal do we try to solve first?
To address the first question, Aesop associates with each rule a \emph{success probability} between 0\% and 100\%.
The success probability is given by the user when registering the rule.
It should roughly approximate the probability that the rule, when successfully applied to an Aesop-provable goal, will lead to a successful proof of that goal.
For example, $∧$-introduction, which reduces a goal $P ∧ Q$ to subgoals $P$ and $Q$, has success probability 100\%: if $P ∧ Q$ is Aesop-provable, then almost certainly $P$ and $Q$ are as well.
Meanwhile, left $∨$-introduction, which reduces a goal $P ∨ Q$ to just $P$, has success probability 50\%, since $P$ and $Q$ are a priori equally likely to be Aesop-provable.
In practice, the exact percentages generally do not affect the proof search much, so we do not worry about whether the empirical success probability of left $∨$-introduction might be 57\% rather than 50\%.

Using the success probabilities, Aesop implements a simple \emph{best-first} strategy, answering the question of which goal to solve first.
Each goal is associated with a \emph{priority} (also a percentage), with the root goal having priority 100\%.
Whenever we apply a rule with success probability $p_r$ to a parent goal with success probability $p_g$, the priority of each subgoal is $p_g p_r$.
Hence, applying a 50\% rule to a 70\% goal results in subgoals with priority 35\%.
Alternatively, rules can assign explicit priorities to the child goals, but this is rarely done.
Depth-first and breadth-first search strategies are likewise supported but little used.

Other Aesop-style tools generally support only depth-first search, which is particularly easy to implement: rule applications are simply backtracked once it becomes clear that they will not lead to a successful proof.
No explicit representation of the search tree is needed; it is encoded in the recursive calls of the search procedure.
ITPs tend to make it easy to take and restore snapshots of the current proof state (and hence to backtrack) since this capability is also used to implement other features, e.g.\ interactive proof development environments.

However, depth-first search scales poorly to larger proofs and rulesets.
If rules are not precisely prioritised, depth-first search may proceed down a long, ultimately fruitless branch, without ever considering a short, successful branch.
Aesop's implementation of best-first search, which is often close to breadth-first in practice, avoids this issue.
Best-first search also allows us to use certain classes of rules that, with a depth-first approach, are either unusable or require very specific configuration.
For example, a transitivity rule such as $a < b → b < c → a < c$ is dangerous because it can be applied infinitely often to its own subgoals.
In Aesop, we can mitigate this danger by giving the rule a very low success probability, so it is tried only as a last resort and the priority of its descendants decreases exponentially with their depth.

Since we support non-depth-first search in Aesop, we must represent the search tree explicitly.
This substantially complicates the handling of \emph{metavariables}.
A metavariable (also called \emph{existential variable}, \emph{schematic variable} and sometimes just \emph{free variable}) is a named hole in the proof state, to be filled by later proof steps.
It often arises from an existential quantifier.
For example, the goal $\Ex{n : ℕ}{0 < n ∧ n < 3}$ is decomposed by built-in Aesop rules into two goals $0 < \mv{n}$ and $\mv{n} < 2$, where $\mv{n}$ is a newly created metavariables.
When a goal is solved, metavariables may be assigned as a side effect.
In our example, applying the lemma $0 < 1$ to the first goal, $0 < \mv{n}$, solves the goal and assigns $\mv{n} ≔ 1$.
The second goal, $\mv{n} < 2$, then becomes $1 < 2$.

This interaction illustrates the fundamental challenge that metavariables pose during tree search: sibling goals are \emph{coupled}, i.e.\ no longer independent.
How the goal $0 < \mv{n}$ is solved determines how its sibling goal $\mv{n} < 2$ is solved, and indeed whether it can be solved at all: we could have applied the fact $0 < 3$ to the first goal, assigning $\mv{n} ≔ 3$ and making the second goal unprovable.

For backtracking-based depth-first search, this is no great obstacle because goals are naturally ordered.
The second goal is proved after the first, taking into account any metavariable assignments made by the first goal's proof.
If the second goal cannot be proved with this assignment, the rules applied to the first goal -- and any metavariable assignments generated by them -- are backtracked and we try again.
Best-first search, however, admits no such fixed goal order.
We may work on the goals $0 < \mv{n}$ and $\mv{n} < 2$ in any order, and we may even interleave the two proof attempts.
Hence, metavariable assignments must be handled explicitly.

Aesop introduces a novel solution for this problem.
When a rule application assigns a metavariable, sibling goals dependent on that metavariable (and certain non-sibling goals on the same search tree branch) are copied and treated as additional children of the rule application.
Intuitively, we view goals with updated metavariable assignments as goals that need to be solved \emph{because} the assignment has been made, and that would not need to be solved if a different assignment had been made.
Hence, consider an attempt to solve our running example with $0 < 3$ applied to the first goal, $0 < \mv{n}$.
This assigns $\mv{n} ≔ 3$, so the coupled sibling goal $\mv{n} < 2$ is copied, with $\mv{n}$ replaced by its assignment.
This means that the rule, which closes its goal and so would ordinarily have no subgoals, produces the subgoal $3 < 2$.
Since that subgoal is unprovable, the attempt will not succeed, but the original goal $0 < \mv{n}$ is still present in the search tree, and Aesop can start a new attempt where this goal is solved by $0 < 1$.
That rule application produces the copied subgoal $1 < 2$ and once this is solved, the overall search succeeds.
We conjecture that this approach is complete, in the sense that if there is a sequence of rule applications that will solve a given goal, a fair search will find it.

As an optimisation, we partition the subgoals of a rule into \emph{metavariable clusters}, which are minimal sets of goals transitively coupled by metavariables.
For example, the goals $0 < \mv{n}$, $\mv{n} < 2$ and $\mv{m} < 3$ are partitioned into two clusters $\{0 < \mv{n}, \mv{n} < 2\}$ and $\{\mv{m} < 3\}$, reflecting the fact that the third goal is independent of the first two.
This increases search parallelism and avoids unnecessary copying.

\section{Efficient Forward Rules}

Since its inception, Aesop has supported \emph{forward rules}.
These are rules that derive facts from the hypotheses available in an Aesop goal and add them to the local context for use by other rules.
For example, given a function $f$, a lemma $\Var{h} : \All{x : ℝ}{0 < f~x}$ and a context containing $a, b : ℝ$, the forward rule for $\Var{h}$ adds hypotheses $\Var{h₁} : 0 < f~a$ and $\Var{h₂} : 0 < f~b$.
These could be used, for example, by Lean's \tac{positivity} tactic to establish positivity of expressions containing $f~a$ and $f~b$.

Aesop initially only provided a very naive implementation of forward rules.
With that implementation, when considering a goal with hypotheses $\Var{h₁} : H₁, \dots, \Var{hₙ} : Hₙ$ during the search, Aesop checks for each forward rule $r : \All{(x₁ : X₁) \dots (xₘ : Xₘ)}{Y}$ whether there is a type-correct substitution $x₁ ↦ \Var{h_{i₁}}, \dots, xₘ ↦ \Var{h_{iₘ}}$, i.e.\ a type-correct assignment of hypotheses to premises of $r$.
If so, the conclusion of $r$, $Y[\Var{h_{i₁}}, \dots, \Var{h_{iₘ}}]$, is added as a new hypothesis (unless it is already present).

A crucial weakness of this scheme is that it is stateless: the matching process is repeated for every goal and every forward rule (subject to an indexing scheme), even when the same rule was already tried on essentially the same goal.
This is a common case: many rules leave the local context entirely unchanged, forward rules themselves only add hypotheses, and even Aesop's \tac{simp} pass usually changes only select hypotheses.

Recent versions of Aesop therefore use a novel incremental mechanism for forward rules.
Each goal is associated with a \emph{forward state} that indexes partial applications of forward rules to the goal's hypotheses.
For instance, given a forward rule $r : \All{x~y : ℕ}{x < y → y < x → ⊥}$ and a goal with hypotheses $a : ℕ,\, b : ℕ,\, \Var{h₁} : a < b$, the forward state stores the fact that $r~a~b~\Var{h₁}$ is a type-correct partial application of $r$.
If a rule then creates a subgoal with additional hypothesis $\Var{h₂} : b < a$, the stored partial application is completed and the conclusion of $r$, $⊥$, is added as a new hypothesis.

To make this process efficient, we index each premise of a forward rule in a discrimination tree~\cite{DiscriminationTrees}.
This allows us to quickly determine that the hypothesis $\Var{h₂} : b < a$ matches both non-dependent premises of $r$, $x < y$ and $y < x$.
Additionally, we store the partial rule applications in a custom data structure that indexes them by the values given to dependent premises, here $x$ and $y$.
This allows us to quickly look up partial applications compatible with the substitution $x ↦ a,\, y ↦ b$ that results from matching $\Var{h₂}$ against the rule's second premise, $y < x$.
The lookup returns $r~a~b~\Var{h₁}$, which we therefore extend to $r~a~b~\Var{h₁}~\Var{h₂}$ and report as a complete forward rule application.
If the rule had additional premises, $r~a~b~\Var{h₁}~\Var{h₂}$ would instead be indexed in the forward state, to be completed by future hypotheses.

To evaluate our scheme, we test it on synthetic and non-synthetic benchmarks.
The synthetic benchmarks show that our scheme substantially outperforms the naive approach in scenarios involving heavy use of forward rules.
For a more practical evaluation, we adapt the benchmark used by lean-auto~\cite{lean-auto} to run different Aesop configurations on Mathlib theorems.
In this benchmark, the new, incremental forward reasoning implementation provides a 1.24x mean speedup over the naive one for Aesop, and a 2.08x mean speedup for the \tac{saturate} tactic that only performs forward reasoning.
The modest mean Aesop speedup may be caused by the fact that most theorems Aesop can solve involve only few applications of forward rules, so the caching provided by the incremental implementation does not have a large effect.
We see larger speedups for invocations with deeper search trees.

\section{Script Generation and Optimisation}

When Aesop finds a proof, or at least makes progress towards one, it applies a series of rules that translate fairly directly into lower-level Lean tactics.
For example, an application of a backward rule $r$ corresponds to $\tac{apply}~r$; an application of a forward rule corresponds to $\tac{have}~\Var{h}~\tac{:=}~r~\dots$.
It is therefore easy (if not entirely trivial) to generate a \emph{tactic script} that can replace a given Aesop call.
Doing so has advantages in two scenarios.

First, a script can replace an Aesop call that closes a goal.
The script typically has much better performance since it performs no search and can optimise certain tactics, e.g.\ use $\tac{simp~only}~[\dots]$ instead of $\tac{simp}$.
It is also stable with respect to changes to the Aesop rule set (though, conversely, less flexible when definitions or lemmas change).

Second, a script for an Aesop call that does not find a full proof may serve as the start of a manual proof.
Hence, Aesop can be used as a goal preprocessor that performs some obvious steps and leaves the rest of the proof to the user.
Doing so without a script is also possible, but creates maintenance challenges since changes to the Aesop rule set are likely to break proofs in ways that are difficult to debug.
Hence, script generation opens up an entirely new use case for Aesop, beyond the closing of straightforward goals that is its traditional remit.

However, if we naively render the used rules as a sequence of tactics, we end up with highly unidiomatic scripts.
They operate on goals in whatever order Aesop used when visiting the goals during the search, rather than the depth-first order that humans overwhelmingly use.
Relatedly, they do not make use of Lean's many facilities for structuring proofs into tree shapes.
And they often use unidiomatic or inefficient tactics.
This jeopardises both use cases since users have to tediously adjust the scripts by hand.

We address these challenges (and more) by adding a script optimiser to Aesop.
It bears some similarity to an optimising compiler or autoformatter and operates in three stages.

First, individual tactics are rewritten into a more performant or more idiomatic form.
For example, we try to use the stock \tac{unfold} tactic instead of Aesop's \tac{aesop\_unfold}, which has slightly different semantics for performance reasons.
Since the semantics are different, we check whether the two tactics actually produce the same goals, which they usually do.
This check is possible because tactics are approximately pure functions of their input goals, and unlike a typical optimising compiler, we know that our script only needs to work for one specific set of initial goals.

Second, we structure the script, reordering the tactics into depth-first order and making use of Lean's structuring facilities, e.g.\ the \tac{case}, \tac{next} and $\tacbullet$ tactic combinators.
Reordering is not always possible when goals contain metavariables, since tactics may behave differently depending on whether a metavariable has already been solved in the proof of another goal.
We propose two methods to address this issue: a static method that misses some opportunities for reordering, and a dynamic method that reorders as much as possible.
The second method is dynamic in the sense that it tries every possible reordering, checking that the reordered tactics have the same behaviour as the original ones.
This makes it more precise, but also much slower.
However, an evaluation on Aesop calls in Mathlib shows that the dynamic method is generally still performant enough for interactive use.

Third, we run two purely syntactic optimisation passes on the resulting script.
These do not refer to the input and output goals of tactics, so are conceptually simple.

\section{Related Work}

Aesop-style provers lie at the intersection of interactive and automated theorem proving, two fields that are almost as old as computing itself.
In the interest of brevity, the following survey of related work is limited to contemporary systems usable from ITPs.
More detailed discussions of work related to Aesop's specific features appear in each of the following chapters.

\subsection{Other Heuristic Provers}

Aesop is most similar to Rocq's and Isabelle's auto tactics, which also perform tree-based search with user-specified rules.
Like Isabelle's auto, Aesop integrates Lean's simplifier, a widely used tactic that performs conditional rewriting with user-specified equations.
Aesop's distinguishing features are those highlighted above: best-first instead of depth-first search with correct and complete metavariable handling; a novel, more scalable treatment of forward rules; and tactic script generation.
Matita's auto~\cite{Matita} augments Rocq's auto with a superposition calculus for equational reasoning and provides a GUI that allows users to inspect the search tree; these extensions could be fruitful for Aesop as well.

Other heuristic, customisable provers include ACL2's \enquote{waterfall}~\cite{ACL2}, PVS's grind~\cite{PVS} and, more distantly related, Mizar's verifier~\cite{MizarOverview2015}.
They serve as the main automation methods of their ITPs and include many extensions that increase their effectiveness.
In particular, their support for arithmetic and equational reasoning is often much better than that of Aesop, which in these domains relies almost entirely on the simplifier.
However, the extensions also add substantial complexity that makes the provers semi-black-box systems in practice.
A notable heuristic prover based not on tree search but saturation is Isabelle's auto2~\cite{auto2}.

\subsection{Automated Theorem Provers and Hammers}

A disadvantage of Aesop and similar tools is that they require extensive, non-trivial customisation to be effective.
By contrast, \emph{automated theorem provers} (ATPs), including SAT and SMT solvers~\cite{SATHandbook,Z3,CVC5} and superposition-based theorem provers~\cite{Vampire,Zipperposition}, provide \enquote{push-button} automation that seeks to find proofs without any help from users.
For certain families of statements, these tools are highly effective, generating complex proofs in seconds.
However, they tend to perform less well on more mathematically sophisticated lemmas.

One way to use an ATP from within an ITP is to implement a \emph{native} ATP in the language of the ITP, providing automated theorem proving as a regular tactic.
Examples of this approach include Agda's~\cite{Agda} Agsy~\cite{Agsy}, HOL4's~\cite{HOL4} Metis~\cite{MetisHOL4} (also available in HOL Light~\cite{HOLLight,MetisHOLLight} and Isabelle~\cite{MetisIsabelle}), Rocq's sauto~\cite{sauto} and Lean's Canonical~\cite{Canonical}, Duper~\cite{Duper} and grind\footnote{\url{https://lean-lang.org/doc/reference/latest/find/?domain=Verso.Genre.Manual.section&name=grind-tactic}}.
However, these ATPs tend to be somewhat weak since they must be re-implemented for each ITP and are therefore less mature than established external ATPs such as CVC5~\cite{CVC5}, Vampire~\cite{Vampire}, Zipperposition~\cite{Zipperposition} or Z3~\cite{Z3}.
Moreover, native ATPs must support their ITPs' logical foundations, rather than a standard logic such as first-order logic (FOL) or higher-order logic (HOL).
This is a particular concern for type-theoretic ITPs, where the foundation is considerably more complex than FOL or HOL.
As a result, native ATPs must adapt established calculi to the type-theoretic setting (Duper, Lean's grind) or rely on simple calculi based on type inhabitation (Agsy, Canonical) that are typically able to find only very short proofs.

The main alternative to native ATPs are external ATPs integrated via \emph{hammers} such as Rocq's CoqHammer~\cite{CoqHammer}, HOL Light's and HOL 4's HOL(y)Hammer~\cite{HOLyHammer,HOLyHammerHOL4}, Lean's LeanHammer~\cite{LeanHammer} and Isabelle's Sledgehammer~\cite{Sledgehammer}.
They prove theorems in four steps.
First, the hammer translates a problem from the logic of the ITP to the (usually weaker) logic of the ATP.
This translation is non-trivial and usually incomplete, but sophisticated heuristics allow for successful translation in many cases.
Second, the hammer heuristically selects a small set of theorems, or \emph{premises}, that are likely to be useful for the proof, and translates them as well.
Third, the hammer asks the ATP to find a proof of the translated problem, given the translated premises.
Fourth, assuming a proof was found, the hammer uses this proof to derive a proof in the ITP.
This so-called \emph{reconstruction} is challenging because ATPs typically do not provide detailed proofs but only the theorems used in their derivations.
Proof reconstruction is then done by passing the theorems to a native ATP, which often succeeds because it enjoys the benefit of very precise premise selection.
Aesop-style provers have sometimes been used as such native ATPs; e.g.\ Rocq's sauto is expressly designed for this purpose, and Aesop is used by LeanHammer as one reconstruction backend among others.
More recently, some ATPs have gained the ability to report detailed proofs~\cite{Alethe,CVC5ProofVerification,CVC5RewriteVerification} that can be directly translated back into the ITP's logic.

Yet another style of ATP integration is exemplified by the Dafny~\cite{Dafny}, \Fstar~\cite{Fstar}, Verus~\cite{Verus} and Why3~\cite{Why3} tools, which are deeply integrated with SMT solvers and allow users to provide only a coarse proof sketch, with the SMT solver filling in most of the details.
The tight integration of ITP and ATP leads to a higher degree of proof automation, but in return, the external ATP is typically trusted, meaning its proofs are not checked by the ITP.
This yields a large and complex trusted code base compared to HOL, Isabelle, Lean and Rocq, all of which strive (to varying degrees) to keep their trusted codebases -- the kernels -- small.

However an ATP is integrated, users pay for the convenience of push-button automation with some amount of brittleness: proofs that compile on one version of the ITP-ATP combination may break or become non-reproducible on the next version, due to updates to the ATP's calculus and heuristics or the hammer's translation algorithms.
Since ATPs employ complex calculi with many heuristics and options, it is often hard to tell what should be done to fix such breakage.
By contrast, proof search procedures like Aesop are much simpler, so it is easier -- though still nontrivial -- to understand why a particular statement could or could not be proved.
Their relative simplicity also makes it easier to obtain consistently good performance, whereas updates to the heuristics of ATPs may influence their performance in complex ways.
This is particularly important in the area of software verification, where proofs are sometimes so large that precise control over the proof search strategy becomes essential.

The flip side of Aesop's simplicity is that, even when properly configured, it is considerably less powerful than ATPs.
For example, Aesop can only make use of equations by passing them to Lean's simplifier, which is helpful only when the equation has a natural left-to-right or right-to-left direction.
By contrast, equality is a central concept in both SMT solvers (via e-matching~\cite{EMatching}) and superposition calculi~\cite{Superposition}.

An evaluation of lean-auto on Mathlib theorems~\cite{lean-auto}, performed by the lean-auto developers, illustrates the tradeoff.
For each of 149,135 Mathlib theorems, several tools are run in an attempt to find a proof.
The tools are given exactly the theorems used by the human-written Mathlib proof, simulating perfect premise selection.
In Aesop's case, this means that the theorems are provided as unsafe backward rules.
With this setup, Aesop, despite not being designed for this use case, proves 47,060 theorems (32\%) in, on average, \SI{93.5}{ms}.
Meanwhile, lean-auto, backed by the Duper, Z3, CVC5 and Zipperposition ATPs, proves 61,906 theorems (42\%), but it also takes an average of \SI{756.8}{ms}.

Lean's new grind tactic (unrelated to PVS's grind) occupies an interesting point in the design space between heuristic provers and classical ATPs.
Like Aesop, it relies on a manually curated set of lemmas tagged with various attributes, giving users the opportunity (and responsibility) to select a suitable tradeoff between power and speed.
However, instead of Aesop's tree-based search, grind uses SMT techniques such as e-matching, heuristic quantifier instantiation and theories for arithmetic, orders, etc.
These techniques make it highly effective at equational and arithmetic reasoning, two weaknesses of Aesop.

\subsection{Neural Theorem Provers}

Recent advances in generative artificial intelligence have led to much interest in \emph{neural theorem provers} (NTPs): tools that search for formal proofs, in Lean or other ITPs, by employing large language models (LLMs).
Various architectures for such NTPs have been proposed, including one-shot NTPs, where the LLM directly generates an entire proof~\cite{DeepSeekProver,GoedelProver,GoedelProverV2}; tactic-by-tactic NTPs, where LLMs generate individual tactics to form a large search tree similar to Aesop's~\cite{Lample22,LeanDojo,ABEL,Aristotle,AlphaProof,DeepSeekProverV15}; decomposition-based systems that ask LLMs to recursively decompose the problem statement into smaller lemmas until the lemmas can be directly proved~\cite{DSP,Hilbert,SeedProver}; and free-form agents that let the LLM control the search process (e.g.\ OpenGauss\footnote{\url{https://www.math.inc/opengauss}} and Leanstral\footnote{\url{https://mistral.ai/news/leanstral}}).
The NTP field is evolving rapidly and many systems are hybrids that combine features of different architectures.

Despite their relative youth, NTPs have already demonstrated capabilities far surpassing those of ATPs.
Notably, several systems have achieved gold and silver medals at the 2024 and 2025 International Mathematical Olympiads.
However, a downside of LLM-based approaches is that they are much more costly, in both time and money, than ATPs.
Typical evaluations of one-shot NTPs use a pass@$k$ metric that gives the system $k$ attempts to solve each problem in a given benchmark, with $k$ ranging from 2 to over 1000.
Even a single such one-shot attempt is typically more expensive than an ATP invocation since LLMs of sufficient size are slow and require costly hardware.
Systems based on tactic-level search or decomposition also tend to involve large numbers of LLM queries.
For instance, the Aristotle system~\cite{Aristotle}, in its IMO gold medal performance, ran Lean on up to 500,000 CPUs for multiple days to verify proofs generated by LLMs running on an undisclosed number of GPUs.
It therefore seems likely that future automated proving systems will incorporate both high-level LLM guidance and efficient ATPs.
The ATPs become, in effect, an optimisation that allows an NTP to shift some of the proving effort from expensive GPUs to cheap CPUs.

\section{Adoption}

For a pragmatic, heuristic tool like Aesop, the ultimate success metric is adoption, so it is my pleasure to report that Aesop has been used for a variety of purposes throughout the Lean ecosystem.
Recent versions of Mathlib call Aesop around 7000 times, often through wrapper tactics that use a specialised rule set.
The \tac{measurability}\footnote{\url{https://github.com/leanprover-community/mathlib4/blob/6b3e5a3312983f468f6187068fca3e69ddc7f80b/Mathlib/Tactic/Measurability.lean}} tactic, which attempts to prove that a function is measurable, is a typical example.
For some such use cases, Aesop has been partially replaced with more efficient specialised tactics such as \tac{fun\_prop}\footnote{\url{https://github.com/leanprover-community/mathlib4/blob/6b3e5a3312983f468f6187068fca3e69ddc7f80b/Mathlib/Tactic/FunProp.lean}}.
Anecdotally, the script generation functionality is used regularly for goals where it would be inopportune to leave Aesop in the final proof script (usually for performance reasons), but since script generation leaves no trace in the final proof, we cannot quantify this.

Even so, many Mathlib theories have not yet been thoroughly annotated to build effective Aesop rule sets, limiting the effectiveness of Aesop even on goals where it could perform well.
Forward rules, in particular, are rarely used in Mathlib.
A tool that suggests good Aesop annotations, perhaps using an LLM to determine whether a lemma should be used forward or backward and what its priority should be, could thus provide substantial value to the Lean community.

Outside of Mathlib, Aesop is used in various mathematical formalisation projects, e.g.\ the Equational Theories project \cite{EquationalTheories} and the formalisation of a proof of the polynomial Freiman-Ruzsa conjecture\footnote{\url{https://github.com/teorth/pfr}}.
Aesop is also used in the nascent-in-Lean software verification domain, for example to dispatch verification conditions in the Loom multi-modal verifier~\cite{Loom}.
Geoffrey Irving's \tac{interval} tactic\footnote{\url{https://github.com/girving/interval}} for interval arithmetic with floating point numbers is largely a wrapper around Aesop with a particular rule set.

As mentioned, Aesop is also used in larger, push-button provers.
LeanHammer~\cite{LeanHammer} uses Aesop as one of its proof reconstruction backends.
Some tactic-by-tactic NTPs use Aesop as a backend providing an AND-OR tree with prioritisation support~\cite{LeanDojo,ABEL}.
And of course, many LLM-based systems use Aesop as a comparatively cheap tactic to dispatch easy goals.

An interesting, unusual use of Aesop is explored in unpublished work by Harry Goldstein et al.%
\footnote{\url{https://github.com/hgoldstein95/palamedes-lean}}
Aesop is used as the backend of a synthesis engine for correct-by-construction generators in the context of property-based testing.
Using a custom set of purely constructive rules, Aesop searches for proofs of the existence of generators with particular correctness properties.
If such a proof is found, an executable generator can be immediately extracted.

\section{Attribution}

This thesis is based on three papers that are reproduced without substantive changes, though I have adjusted the introductions, centralised some of the background material and related work and made various small updates.
Chapter~\ref{chap:aesop} is based on \emph{Aesop: White-Box Best-First Proof Search for Lean}~\cite{Aesop}, which is joint work with Asta Halkjær From.
I have added Sec.~\ref{sec:aesop:mvars-summary} to give a more precise description of our handling of metavariables.
Chapter~\ref{chap:forward} is based on \emph{Incremental Forward Reasoning for White-Box Proof Search}~\cite{Aesop-forward}, which is joint work with Xavier Généreux.
I have added Sec.~\ref{sec:forward:rpinf}, which discusses an attempt at performance optimisation that did not make it into the final implementation.
This section was intended to be part of the paper, but was removed due to space constraints.
Chapter~\ref{chap:script} is based on \emph{Tactic Script Optimisation for Aesop}~\cite{Aesop-script}.
For consistency, I use the first person plural throughout, but any mistakes in this thesis are, of course, solely my own.
I used AI tools for some editing tasks and to write data analysis scripts for the evaluation of Chapter~\ref{chap:forward}.

\section{Data Availability}

All functionality described in this thesis is implemented in mainline Aesop\footnote{\url{https://github.com/leanprover-community/aesop}}.
The exact Aesop version corresponding to each chapter, as well as benchmark code and data, are available in the supplement for the paper that the respective chapter is based on~\cite{AesopSupplement,ForwardSupplement,ScriptSupplement}.
